\lyseis{3}{Δυνάμεις και προτεραιότητα: ποιος πάει πρώτος}
\sol{1} α) $5^3$ \quad β) $2^5$ \quad γ) $10^4$ \quad δ) $7^2$ \quad ε) $9=9^1$ (ή $3^2$).
\sol{2} α) $2\cdot 2\cdot 2\cdot 2=16$ \quad β) $3\cdot 3\cdot 3=27$ \quad γ) $5\cdot 5=25$ \quad δ) $100.000$ \quad ε) $1$ \quad στ) $6$.
\sol{3} Στο $4^5$: βάση 4, εκθέτης 5. Στο $12^2$: βάση 12, εκθέτης 2. Στο $1^9$: βάση 1, εκθέτης 9. Στο $10^8$: βάση 10, εκθέτης 8.
\sol{4} Για τους 2, 3, 4, 5, 10: τετράγωνα 4, 9, 16, 25, 100 και κύβοι 8, 27, 64, 125, 1.000.
\sol{5} α) $12+12=24$ \quad β) $15\cdot 4=60$ \quad γ) $24:4=6$, μετά $6:2=3$ \quad δ) $100-20=80$ \quad ε) $2+9=11$ \quad στ) $5^2=25$.
\sol{6} α) $10.100$ (δες το παράδειγμα). β) $5^3-8\cdot 9=125-72=53$. γ) $3\cdot 6^2-10\cdot 3=3\cdot 36-30=108-30=78$.
\sol{7} α) $7\cdot 10^3+6\cdot 10^2+0\cdot 10+4$ \quad β) $5\cdot 10^4+3\cdot 10^3+2\cdot 10^2+8\cdot 10$ \quad γ) $4\cdot 10^5+6\cdot 10+5$.
\sol{8} α) $3.205$ \quad β) $900.040$ \quad γ) $1.000.100$.
\sol{9} α) Λάθος: $2^3=8$. β) Σωστό. γ) Λάθος: $5^2=25$, ενώ $2\cdot 5=10$. δ) Σωστό. ε) Λάθος: πρώτα ο πολλαπλασιασμός, $3+10=13$. στ) Λάθος: $(1+1)^2=4$, ενώ $1^2+1^2=2$.
\sol{10} Με 3 τροχούς: $10^3=1.000$ κωδικοί. Με 4: $10^4=10.000$. Με 5: $10^5=100.000$. Κάθε τροχός που προστίθεται πολλαπλασιάζει τους κωδικούς επί 10.
\sol{11} Με 3 διακόπτες: $2^3=8$ συνδυασμοί. Με 5: $2^5=32$. Με 10: $2^{10}=1.024$. Κάθε διακόπτης που προστίθεται διπλασιάζει τους συνδυασμούς.
\sol{12} $4^3=4\cdot 4\cdot 4=64$ κυβικά δεκατόμετρα, δηλαδή 64 λίτρα.
\sol{13} Εμβαδό $12^2=144$ τ.μ. Με πλευρά 24 μ. το εμβαδό είναι $24^2=576$ τ.μ., που είναι $4\cdot 144$. Όχι διπλάσιο, τετραπλάσιο.
\sol{14} $8^2=64$ τετράγωνα.
\sol{15} Ταιριάζει η $3\cdot 8+12$: τρεις μπάλες των 8 € και ένα φανελάκι των 12 €. Πληρώνει $24+12=36$ €. Η $3\cdot(8+12)=60$ € θα ταίριαζε αν αγόραζε τρία ζευγάρια «μπάλα και φανελάκι».
\sol{16} Σωστή είναι η Ζωή (300). Στην παράσταση $4\cdot(7+7\cdot 9)+20$ πρώτα γίνεται ο πολλαπλασιασμός μέσα στην παρένθεση: $7\cdot 9=63$, μετά $7+63=70$, μετά $4\cdot 70=280$ και $280+20=300$. Ο Ηλίας (524) έκανε πρώτα το $7+7=14$ μέσα στην παρένθεση: $4\cdot 14\cdot 9+20=524$. Ο Άρης (335) έκανε πρώτα το $4\cdot 7+7=35$ και μετά επί 9: $35\cdot 9+20=335$. Και οι δύο παραβίασαν την προτεραιότητα.
\sol{17} $2^3\cdot 2^2=8\cdot 4=32$. Ως παράγοντες: $(2\cdot 2\cdot 2)\cdot(2\cdot 2)$, δηλαδή πέντε δυάρια: $2^5=32$. Ο κανόνας: προσθέτω τους εκθέτες, $2^3\cdot 2^2=2^{3+2}$.
\sol{18} $2^5=32>5^2=25$ και $2^4=16=4^2$. Ο $\alpha=0$ κάνει το $\alpha^2=0$ ίσο με το $2\alpha=0$. (Μόνο το 0 και το 2 ικανοποιούν την ισότητα.)
\sol{19} Χωρίς παρενθέσεις: $2+12+5=19$. Με παρενθέσεις: $(2+3)\cdot 4+5=25$, $2+3\cdot(4+5)=2+27=29$, $(2+3)\cdot(4+5)=45$. Οι διαφορετικές τιμές είναι 19, 25, 29 και 45.
\sol{20} Π.χ. $1+2+3+4=10$ και $1\cdot 2\cdot 3+4=10$.
\sol{21} Το μοτίβο: $10^3=1.000$, $10^2=100$, $10^1=10$. Κάθε φορά που κατεβαίνει ο εκθέτης κατά 1, ο αριθμός διαιρείται με το 10. Το επόμενο είναι $10:10=1$, άρα $10^0=1$. Δίκιο έχει η Ζωή.
